% Bagian: second-order-logic
% Bab: syntax-and-semantics
% Subbagian: semantic-notions

\documentclass[../../../include/open-logic-section]{subfiles}

\begin{document}

\olfileid[id]{sol}{syn}{sem}
\olsection{Gagasan Semantis}

\begin{explain}
Gagasan logis utama berupa \emph{validitas}, \emph{perikutan}, dan
\emph{keterpenuhan} didefinisikan dengan cara yang sama bagi logika orde kedua
seperti bagi logika orde pertama, kecuali bahwa relasi keterpenuhan yang
mendasarinya sekarang adalah relasi bagi !!{formula}s orde kedua. Tentu saja,
suatu !!{sentence} orde kedua merupakan !!a{formula} yang semua variabelnya,
termasuk variabel predikat dan fungsi, terikat.
\end{explain}

\begin{defn}[Validitas]
Suatu kalimat $!A$ \emph{valid}, $\Entails !A$, jika dan hanya jika
$\Sat{M}{!A}$ bagi setiap !!{structure}~$\Struct M$.
\end{defn}

\begin{defn}[Perikutan]
Suatu himpunan kalimat~$\Gamma$ \emph{mengakibatkan} suatu kalimat~$!A$,
$\Gamma \Entails !A$, jika dan hanya jika bagi setiap
!!{structure}~$\Struct M$ dengan $\Sat{M}{\Gamma}$, berlaku
$\Sat{M}{!A}$.
\end{defn}

\begin{defn}[Keterpenuhan]
Suatu himpunan kalimat~$\Gamma$ \emph{dapat dipenuhi} jika
$\Sat{M}{\Gamma}$ bagi suatu !!{structure}~$\Struct M$. Jika $\Gamma$ tidak
dapat dipenuhi, himpunan itu disebut \emph{tak dapat dipenuhi}.
\end{defn}

\end{document}
